\subsection{দ্বৈত রূপ, পূরক ও সরলীকরণের যুক্তি}
দ্বৈত রূপ বের করতে অপারেশন ও ধ্রুবক বদলায়, চলকের ওপর নিজের থেকে পূরক বসে না। রাশির পূরক নির্ণয়ে সম্পূর্ণ রাশির মান ০ থেকে ১ বা ১ থেকে ০ হয়। এই দুই প্রক্রিয়া পৃথক।

\chThreeExample{১}{দ্বৈত রূপ এবং পূরক}
$F=A+BC$-এর দ্বৈত রূপ ও পূরক বের করো।
\chThreeAnswer দ্বৈত রূপ $A(B+C)$। ডি-মরগ্যান প্রয়োগে পূরক $\overline F=\overline A(\overline B+\overline C)$। দ্বৈত রূপ এবং মূল রাশির মান সব ইনপুটে একই হবে—এমন দাবি নেই; বৈধ \emph{সূত্রের} দ্বৈত সূত্রও বৈধ।

\smallhead{প্রয়োজনীয় সরলীকরণ সূত্রের কারণ}
\[\begin{aligned}
A+\overline AB&=(A+\overline A)(A+B)=A+B,\\
(A+B)(A+\overline B)&=A+B\overline B=A.
\end{aligned}\]
এখানে প্রথম লাইনে $A+BC=(A+B)(A+C)$, দ্বিতীয় লাইনে তার বিপরীত ব্যবহার হয়েছে। সাধারণ বীজগণিতের মতো $1+1=2$ ধরে হিসাব করলে ভুল হবে।

\smallhead{Consensus সূত্র}
\[AB+\overline AC+BC=AB+\overline AC.\]
কারণ $BC=BC(A+\overline A)=ABC+\overline ABC$। প্রথম নতুন পদ $AB$-এর এবং দ্বিতীয়টি $\overline AC$-এর মধ্যে আগেই অন্তর্ভুক্ত। ফলে $BC$ নতুন কোনো ইনপুটে আউটপুট ১ করছে না। দ্বৈত সূত্র $(A+B)(\overline A+C)(B+C)=(A+B)(\overline A+C)$।

\chThreeExample{২}{সরলীকরণে বাস্তব সিদ্ধান্ত অপরিবর্তিত}
এক যন্ত্র চালু হয় যদি প্রধান সুইচ $A$ চালু থাকে, অথবা $A$ ও অনুমতি সুইচ $B$ দুটোই চালু থাকে।
\chThreeAnswer $F=A+AB=A(1+B)=A$। দ্বিতীয় শর্ত প্রথমটির মধ্যেই অন্তর্ভুক্ত; $B$-এর জন্য আলাদা গেট দরকার নেই। সরলীকরণে শর্তের ফল বদলায়নি।

\chThreeExample{৩}{ডি-মরগ্যানে পূরকের সীমা}
$F=\overline{A+\overline{BC}}$ সরল করো।
\chThreeAnswer $F=\overline A\,\overline{\overline{BC}}=\overline ABC$। ভেতরের পুরো $BC$-এর দ্বৈত পূরক উঠে যায়; শুধু B-এর পূরক তোলা যাবে না।

\chThreeExample{৪}{ধারণাভিত্তিক MCQ}
বুলিয়ান বণ্টন সূত্র কোনটি? (ক) $A+BC=(A+B)(A+C)$, (খ) $A+A=2A$, (গ) $A+1=A$, (ঘ) $A\overline A=1$।
\chThreeAnswer (ক)। (খ), (গ), (ঘ)-তে যথাক্রমে $A$, $1$ ও $0$ সঠিক ফল।
